\subsection{Main Rock Types Relevant to Thermoelastic Wave Propagation}

Among the many rock types found in the crust, sandstone, basalt, granite, and limestone provide a useful set for introductory analysis because they represent distinct geological origins and contrasting physical architectures. Sandstone is typically clastic and granular, basalt is volcanic and commonly fine-grained, granite is intrusive and crystalline, and limestone is carbonate and often strongly modified by biochemical and diagenetic processes. Together these rocks span a wide range of porosity, pore structure, mineral composition, stiffness, thermal conductivity, and anisotropy. This makes them suitable reference materials for discussing how thermal disturbances and elastic waves may propagate differently through geological media \cite{mavko2009,robertson1988,earle2019}.

Sandstone is a sedimentary rock formed by the deposition, burial, and lithification of sand-sized particles. In many common sandstones, quartz is the dominant framework mineral, with feldspar and lithic fragments present in varying amounts. The grains are bound by cement, which may be silica, carbonate, clay, or iron oxide. The framework is commonly grain-supported, so the mechanical role of grain contacts is especially important. Sandstone is therefore a classic example of a medium in which porosity can be significant and in which pore geometry, cementation, and fluid saturation strongly influence effective stiffness and wave velocity. Clean quartz arenites, feldspathic sandstones, and clay-rich or carbonate-cemented sandstones may all behave differently even when grouped under the same general rock name. Diagenesis is particularly important: quartz overgrowths, carbonate cementation, clay authigenesis, and dissolution can either destroy or create porosity and permeability. Because of bedding, sorting variations, compaction fabrics, and deformation bands in porous granular rocks, sandstone often shows both heterogeneity and directional behavior. In thermoelastic problems, quartz-rich composition also matters because quartz-bearing rocks commonly have distinctive thermal conductivity behavior and a noticeable sensitivity of bulk properties to porosity and saturation \cite{earle2019,mavko2009,robertson1988}.

Basalt is an igneous volcanic rock produced by the cooling of mafic lava at or near the Earth's surface. It is usually dark, fine-grained, and dominated by plagioclase and pyroxene, with olivine and iron-titanium oxides commonly present in many varieties. In comparison with many sandstones, intact basalt is typically denser and less porous at the matrix scale, although volcanic textures can add complexity through vesicles, interflow zones, and alteration products. Basaltic lava flows also develop cooling fractures, and where cooling occurs under suitable conditions, contraction produces columnar jointing. These joints and cooling-related fracture networks are especially important for directional behavior because they introduce preferred pathways and preferred weaknesses into a rock that might otherwise be treated as relatively massive. From a thermoelastic perspective, basalt is important as a dense mafic rock whose thermal and mechanical response differs from quartz-rich sedimentary or felsic crystalline rocks. Basalt may combine relatively high stiffness of the intact matrix with strong fracture control at the flow scale, which makes its effective behavior dependent on both matrix properties and cooling-related structural features \cite{earle2019,robertson1988,turcotte2014}.

Granite is an intrusive igneous rock formed by the slow cooling and crystallization of silica-rich magma at depth. It has a crystalline, coarse-grained texture and is composed mainly of quartz, feldspar, and mica, with plagioclase and accessory minerals present in variable proportions. Because the crystals are interlocking and the rock forms under plutonic conditions, intact granite generally has low primary porosity and high stiffness compared with many porous sedimentary rocks. At the rock-mass scale, however, fluid flow and mechanical weakness are commonly controlled by joints, fractures, exfoliation planes, and weathered zones rather than by matrix pores. Granite is especially relevant to thermoelastic analysis because it is polymineralic: quartz, feldspars, and micas do not expand thermally in identical ways. Differences in thermal expansion among neighboring minerals can generate internal stresses and promote microcracking, particularly under repeated or strong temperature changes. Accordingly, granite may appear mechanically competent in intact form, yet develop internal thermal damage due to mineral-property mismatch and external heterogeneity due to fractures and weathering. These features make granite a useful reference medium for studying the coupling of thermal expansion, stress concentration, and elastic-wave response in crystalline rocks \cite{earle2019,jaeger2007,robertson1988}.

Limestone is a carbonate sedimentary rock composed chiefly of calcite and, in some cases, aragonite-derived carbonate material. It may form as a biochemical rock from shell and skeletal debris, as a chemically precipitated rock, or through combinations of depositional and early diagenetic processes. This implies wide textural diversity: some limestones are compact and fine-grained, whereas others are fossiliferous, oolitic, fractured, chalky, recrystallized, or strongly vuggy. Carbonate rocks differ fundamentally from siliciclastic rocks because their pore systems are often governed less by simple depositional grain packing and more by diagenesis, dissolution, cementation, recrystallization, and stylolitization. As a consequence, their porosity and elastic behavior may vary sharply even among limestones of similar age or depth. Rock-physics treatments of carbonates emphasize that pore type is one of the main controls on acoustic velocity and that carbonate properties are often strongly influenced by diagenetic processes rather than by a simple monotonic burial trend. In thermoelastic terms, limestone is also important because carbonate minerals have their own thermal characteristics, and calcitic rocks may exhibit substantial variability in thermal conductivity and thermal stress response depending on purity, pore content, saturation, and internal structure. Dissolution cavities, fractures, and vugs can further enhance heterogeneity and make both thermal and elastic propagation highly irregular \cite{moore2001,mavko2009,robertson1988}.

Taken together, these four rock types illustrate why thermoelastic wave propagation in geological media cannot be reduced to a single ``rock behavior.'' Sandstone highlights the role of grain-supported porosity, cementation, and saturation. Basalt shows how a dense volcanic matrix can coexist with cooling joints and flow-scale directional structure. Granite emphasizes crystalline stiffness together with grain-scale thermal mismatch and fracture-controlled bulk behavior. Limestone demonstrates the importance of depositional fabric, dissolution, and diagenesis in carbonate systems. In broad comparison, sandstone and limestone are commonly more sensitive to pore structure, pore fluids, and diagenetic modification, whereas basalt and granite are often controlled more strongly by crystalline matrix stiffness and fracture networks. This contrast is important because thermoelastic response is controlled not only by mineral composition, but also by the internal architecture of the rock.

In all cases, density and elastic stiffness influence wave speed, while thermal expansion, thermal conductivity, and heat capacity control how temperature disturbances diffuse and how thermal strain accumulates. Directional behavior may emerge from bedding in sandstone, cooling-joint geometry in basalt, joint and fabric organization in granite, and bedding-fracture-vug networks in limestone. The same thermal loading can therefore produce different stress fields and different patterns of wave propagation depending on lithology and internal structure \cite{mavko2009,robertson1988,aki2002,schultz2019}.

After defining geological media in general terms and introducing these representative rock types, the next section can examine the physical parameters that control thermoelastic wave propagation more directly. In particular, attention must turn to density, porosity, elastic stiffness, Poisson's ratio, thermal expansion coefficient, thermal conductivity, heat capacity, and thermal diffusivity, since it is through these properties that geological origin and rock structure enter the classical theory of thermoelastic response and elastic-wave propagation \cite{aki2002,mavko2009,robertson1988}.